\section*{Acknowledgments}

We would like to thank Lucius Bushnaq and Yannick Kluth for their help
in checking the Feynman rules and the general setup, Jonas Klappert for
his support in using \texttt{FireFly}, Johann Usovitsch for his support
in using \texttt{Kira}, Alexander Voigt for general technical
discussions, and Mauro Papinutto for pointing out a typo in an earlier
version of this manuscript.  This work was supported by \textit{Deutsche
  Forschungsgemeinschaft (DFG)} through project
\href{http://gepris.dfg.de/gepris/projekt/386986591}{HA 2990/9-1}, and
by the U.S.\ Department of Energy under award No.\ DE-SC0008347.  This
document was prepared using the resources of the Fermi National
Accelerator Laboratory (Fermilab), a U.S. Department of Energy, Office
of Science, HEP User Facility. Fermilab is managed by Fermi Research
Alliance, LLC (FRA), acting under Contract No.\ DE-AC02-07CH11359.  Some
of the calculations were performed with computing resources granted by
\abbrev{RWTH} Aachen University under project \textit{rwth0244}.  All
Feynman diagrams in this paper were produced using
Ti\textit{k}Z-Feynman\,\cite{Ellis:2016jkw}.



\appendix

\section{Feynman rules}
\label{sec:frules}

In this section we present all the Feynman rules required for our
calculation.  All momenta are considered to be outgoing.  $\xi$ is the
usual \qcd\ gauge parameter and $\kappa$ the additional gradient-flow
gauge parameter introduced in the flow equations \noeqn{eq:flow}.  Lines
with a solid arrow are flow lines and lines with a dashed arrow can be
both flow lines or flowed propagators. The symbols $p,q,r$ denote
Euclidean four-momenta, while $s,t$ are real-valued flow-time variables.

\textbf{Propagators and flow lines}
\begin{itemize}
\item gluon propagator:
\begin{flalign}
  & \raisebox{-0.5em}{\includegraphics{images/gluon-propagator.pdf}}
  & = &
  \begin{split}
    \delta^{ab} \frac{1}{p^2} \bigg( & \left( \delta_{\mu \nu} - \frac{p_\mu p_\nu}{p^2} \right) e^{-(t+s)p^2} \\
    & + \xi \frac{p_\mu p_\nu}{p^2} e^{-\kappa(t+s)p^2} \bigg)
  \end{split}
  & &
\end{flalign}

\item (anti)quark propagator:
\begin{flalign}
  & \raisebox{-0.5em}{\includegraphics{images/quark-propagator.pdf}}
  & = &
  \begin{split}
    \delta_{ij} \frac{(-\mathrm{i}\slashed{p}+m)_{\alpha\beta}}{p^2+m^2} \, e^{-(t+s)p^2}
  \end{split}
  & &
\end{flalign}

\item ghost propagator:
\begin{flalign}
  & \raisebox{-0.5em}{\includegraphics{images/ghost-propagator.pdf}}
  & = &
  \begin{split}
    \delta^{ab} \frac{1}{p^2}
  \end{split}
  & &
\end{flalign}

\item gluon flow line:
\begin{flalign}
  & \raisebox{-1.2em}{\includegraphics{images/gluon-flowline.pdf}}
  & = &
  \begin{split}
    \delta^{ab} \, \theta(t-s) \bigg( & \left( \delta_{\mu \nu} - \frac{p_\mu p_\nu}{p^2} \right) e^{-(t-s)p^2} \\
    &+ \frac{p_\mu p_\nu}{p^2} \, e^{-\kappa(t-s)p^2} \bigg)
  \end{split}
  & &
\end{flalign}

\item quark flow line:
\begin{flalign}
  & \raisebox{-1.2em}{\includegraphics{images/quark-flowline.pdf}}
  & = &
  \begin{split}
    \delta^{ij} \delta_{\alpha\beta} \, \theta(t-s) \, e^{-(t-s)p^2}
  \end{split}
  & &
\end{flalign}

\item antiquark flow line:
\begin{flalign}
  & \raisebox{-1.2em}{\includegraphics{images/antiquark-flowline.pdf}}
  & = &
  \begin{split}
    \delta^{ij} \delta_{\alpha\beta} \, \theta(t-s) \, e^{-(t-s)p^2}
  \end{split}
  & &
\end{flalign}
\end{itemize}

\textbf{Vertices}
\begin{itemize}
\item three-gluon vertex:
\begin{flalign}
  &
  \begin{gathered}
    \includegraphics{images/ggg-vertex.pdf}
  \end{gathered}
  & = &
  \begin{split}
    \mathrm{i} g f^{abc} \, \big( \delta_{\mu \nu} (p-q)_\rho + \delta_{\nu\rho} (q-r)_\mu
    + \delta_{\rho\mu} (r-p)_\nu \big)
  \end{split}
  & &
  \label{fr:3gluon}
\end{flalign}

\item four-gluon vertex:
\begin{flalign}
  &
  \begin{gathered}
    \includegraphics{images/gggg-vertex.pdf}
  \end{gathered}
  & = &
  \begin{split}
    -g^2 \, \big( & f^{abe} f^{cde} (\delta_{\mu\rho}\delta_{\nu\sigma} - \delta_{\mu\sigma}\delta_{\nu\rho} ) \\
    & + f^{ace} f^{bde} (\delta_{\mu\nu}\delta_{\rho\sigma} - \delta_{\mu\sigma}\delta_{\nu\rho} ) \\
    & + f^{ade} f^{bce} (\delta_{\mu\nu}\delta_{\rho\sigma} - \delta_{\mu\rho}\delta_{\nu\sigma} ) \big)
  \end{split}
  & &
\label{fr:4gluon}
\end{flalign}

\item ghost-gluon vertex:
\begin{flalign}
  &
  \begin{gathered}
    \includegraphics{images/ghost-g-vertex.pdf}
  \end{gathered}
  & = &
  \begin{split}
    - \mathrm{i} g f^{abc} \, p_\mu
  \end{split}
  & &
\end{flalign}

\item quark-gluon vertex:
\begin{flalign}
  &
  \begin{gathered}
    \includegraphics{images/qqg-vertex.pdf}
  \end{gathered}
  & = &
  \begin{split}
    g \, (T^a)_{ij} \, \big( \gamma_\mu \big)_{\alpha\beta}
  \end{split}
  & &
\end{flalign}

\item two-plus-one gluon flow vertex:
\begin{flalign}
  &
  \begin{gathered}
    \includegraphics{images/ggg-flow-vertex.pdf}
  \end{gathered}
  & = &
  \begin{split}
    -\mathrm{i} g f^{abc} \int_0^\infty \mathrm{d}s \, \big( & \delta_{\nu\rho} (r-q)_\mu + 2 \delta_{\mu\nu} q_\rho - 2 \delta_{\mu\rho} r_\nu \\
    & + (\kappa-1) (\delta_{\mu\rho} q_\nu  - \delta_{\mu\nu} r_\rho) \big) \\
  \end{split}
  & &
\end{flalign}

\item three-plus-one gluon flow vertex:
\begin{flalign}
  &
  \begin{gathered}
    \includegraphics{images/gggg-flow-vertex.pdf}
  \end{gathered}
  & = &
  \begin{split}
    -g^2 \int_0^\infty \mathrm{d}s \, \big( & f^{abe} f^{cde} (\delta_{\mu\rho}\delta_{\nu\sigma} - \delta_{\mu\sigma}\delta_{\nu\rho} ) \\
    & + f^{ace} f^{bde} (\delta_{\mu\nu}\delta_{\rho\sigma} - \delta_{\mu\sigma}\delta_{\nu\rho} ) \\
    & + f^{ade} f^{bce} (\delta_{\mu\nu}\delta_{\rho\sigma} - \delta_{\mu\rho}\delta_{\nu\sigma} ) \big)
  \end{split}
  & &
\end{flalign}

\item quark-one-gluon flow vertex:
\begin{flalign}
  &
  \begin{gathered}
    \includegraphics{images/qqg-flow-vertex.pdf}
  \end{gathered}
  & = &
  \begin{split}
    \mathrm{i} g \, \delta_{\alpha\beta} \big( T^a \big)_{ij} \int_0^\infty \mathrm{d} s \, \big( 2 p_\mu + (1-\kappa) q_\mu \big)
  \end{split}
  & &
\end{flalign}

\item antiquark-one-gluon flow vertex:
\begin{flalign}
  &
  \begin{gathered}
    \includegraphics{images/anti-qqg-flow-vertex.pdf}
  \end{gathered}
  & = &
  \begin{split}
    - \mathrm{i} g \, \delta_{\alpha\beta} \big( T^a \big)_{ji} \int_0^\infty \mathrm{d} s \, \big( 2 p_\mu + (1-\kappa) q_\mu \big)
  \end{split}
  & &
\end{flalign}

\item quark-two-gluons flow vertex:
\begin{flalign}
  &
  \begin{gathered}
    \includegraphics{images/qqgg-flow-vertex.pdf}
  \end{gathered}
  & = &
  \begin{split}
    g^2 \delta_{\alpha\beta} \delta_{\mu\nu} \left\{ T^a,T^b \right\}_{ij} \int_0^\infty \mathrm{d}s
  \end{split}
  & &
\end{flalign}

\item antiquark-two-gluons flow vertex:
\begin{flalign}
  &
  \begin{gathered}
    \includegraphics{images/anti-qqgg-flow-vertex.pdf}
  \end{gathered}
  & = &
  \begin{split}
    g^2 \delta_{\alpha\beta} \delta_{\mu\nu} \left\{ T^a,T^b \right\}_{ji} \int_0^\infty \mathrm{d}s
  \end{split}
  & &
\end{flalign}
\end{itemize}

Feynman rules for the composite flow-time operators:
\begin{itemize}
\item $G_{\mu\nu}^a(t,x)G_{\mu\nu}^a(t,x)$:
  \begin{itemize}
  \item \two-gluon vertex:
    \begin{flalign}
      &
      \begin{gathered}
        \includegraphics{images/E_2.pdf}
      \end{gathered}
      & = &
      \begin{split}
        - g^2 \delta^{ab} \big(\delta_{\mu\nu} p \cdot q - p_\mu q_\nu\big)
      \end{split}
      & &
      \label{fr:egg}
    \end{flalign}
  \item \three-gluon vertex:
    \begin{flalign}
      &
      \begin{gathered}
        \includegraphics{images/E_3.pdf}
      \end{gathered}
      & = &
      \begin{split}
        \mathrm{i} g^3 f^{abc} \big( & \delta_{\mu\nu} (q-p)_\rho + \delta_{\nu\rho} (k-q)_\mu \\
        & + \delta_{\mu\rho} (p-k)_\nu \big)
      \end{split}
      & &
      \label{fr:eggg}
    \end{flalign}
  \item \four-gluon vertex:
    \begin{flalign}
      &
      \begin{gathered}
        \includegraphics{images/E_4.pdf}
      \end{gathered}
      & = &
      \begin{split}
        g^4 \big( & f^{abe} f^{cde} (\delta_{\mu\rho}\delta_{\nu\sigma} - \delta_{\mu\sigma}\delta_{\nu\rho} ) \\
        & + f^{ace} f^{bde} (\delta_{\mu\nu}\delta_{\rho\sigma} - \delta_{\mu\sigma}\delta_{\nu\rho} ) \\
        & + f^{ade} f^{bce} (\delta_{\mu\nu}\delta_{\rho\sigma} - \delta_{\mu\rho}\delta_{\nu\sigma} ) \big)
      \end{split}
      & &
      \label{fr:egggg}
    \end{flalign}
  \end{itemize}
\end{itemize}

\begin{itemize}
\item $\bar\chi(t,x)\chi(t,x)$:
    \begin{flalign}
      &
      \begin{gathered}
        \includegraphics{images/S_2.pdf}
      \end{gathered}
      & = &
      \begin{split}
        \delta_{ij} \delta_{\alpha\beta}
      \end{split}
      & &
      \label{fr:sqq}
    \end{flalign}
\end{itemize}

\begin{itemize}
  \item $\bar\chi(t,x)\overleftrightarrow{\mathcal{D}}^\mathrm{F}\chi(t,x)$:
    \begin{itemize}
    \item quark-antiquark vertex:
    \begin{flalign}
      &
      \begin{gathered}
        \includegraphics{images/K_2.pdf}
      \end{gathered}
      & = &
      \begin{split}
        \mathrm{i} \delta^{ij} (\slashed p - \slashed q)_{\alpha\beta}
      \end{split}
      & &
      \label{fr:rqq}
    \end{flalign}
    \item quark-antiquark-gluon vertex:
    \begin{flalign}
      &
      \begin{gathered}
        \includegraphics{images/K_3.pdf}
      \end{gathered}
      & = &
      \begin{split}
        -2 g (T^a)^{ij} (\gamma_\mu)_{\alpha\beta}
      \end{split}
      & &
      \label{fr:rqqg}
    \end{flalign}
    \end{itemize}
\end{itemize}

\textbf{Auxiliary Feynman rules}\\
The following Feynman rules are auxiliary Feynman rules for our implementation as described in \sct{sec:implementation}.
They represent the vertices with four gluons or gluon-flow lines and allow us to factorise the color structure of all Feynman diagrams.

\begin{itemize}
  \item four-gluon vertex:
  \begin{itemize}
    \item $\Sigma$ propagator:
      \begin{flalign}
        &
        \begin{gathered}
          \includegraphics{images/sigma_propagator.pdf}
        \end{gathered}
        & = &
        \begin{split}
          \delta^{ab}\delta_{\mu\rho}\delta_{\nu\sigma}
        \end{split}
        & &
      \end{flalign}
    \item $\Sigma$-gluon-gluon vertex:
      \begin{flalign}
        &
        \begin{gathered}
          \includegraphics{images/sigma_vertex.pdf}
        \end{gathered}
        & = &
        \begin{split}
          \frac{\mathrm{i}g}{\sqrt{2}} f^{abc}\left(\delta_{\mu\rho}\delta_{\nu\sigma}-\delta_{\nu\rho}\delta_{\mu\sigma}\right)
        \end{split}
        & &
      \end{flalign}
  \end{itemize}
  \item three-plus-one gluon flow vertex:
    \begin{itemize}
      \item $\Sigma$ flow line:
        \begin{flalign}
          &
          \begin{gathered}
            \includegraphics{images/sigma_flow_line.pdf}
          \end{gathered}
          & = &
          \begin{split}
            \delta^{ab}\delta_{\mu\rho}\delta_{\nu\sigma} \delta(t-s)
          \end{split}
          & &
        \end{flalign}
      \item outgoing $\Sigma$-gluon-gluon flow vertex:
        \begin{flalign}
          &
          \begin{gathered}
            \includegraphics{images/sigma_flow_vertex_out.pdf}
          \end{gathered}
          & = &
          \begin{split}
            \frac{\mathrm{i}g}{\sqrt{2}} f^{abc}\left(\delta_{\mu\rho}\delta_{\nu\sigma} -\delta_{\nu\rho}\delta_{\mu\sigma}\right) \int_0^\infty \mathrm{d}s
          \end{split}
          & &
        \end{flalign}
      \item incoming $\Sigma$-gluon-gluon flow vertex:
        \begin{flalign}
          &
          \begin{gathered}
            \includegraphics{images/sigma_flow_vertex_in.pdf}
          \end{gathered}
          & = &
          \begin{split}
            \frac{\mathrm{i}g}{\sqrt{2}} f^{abc}\left(\delta_{\mu\rho}\delta_{\nu\sigma}-\delta_{\nu\rho}\delta_{\mu\sigma}\right) \int_0^\infty \mathrm{d}s
          \end{split}
          & &
        \end{flalign}
    \end{itemize}
    \item \four-gluon vertex from $G_{\mu\nu}^a(t,x)G_{\mu\nu}^a(t,x)$:
      \begin{itemize}
        \item $\Sigma_E$ flow line:
          \begin{flalign}
            &
            \begin{gathered}
              \includegraphics{images/E_sigma.pdf}
            \end{gathered}
            & = &
            \begin{split}
              \delta^{ab}\delta_{\mu\rho}\delta_{\nu\sigma} \delta(t-s)
            \end{split}
            & &
          \end{flalign}
        \item outgoing $\Sigma_E$-gluon-gluon vertex:
          \begin{flalign}
            &
            \begin{gathered}
              \includegraphics{images/E_sigma_vertex_out.pdf}
            \end{gathered}
            & = &
            \begin{split}
              \frac{g^2}{\sqrt{2}} f^{abc}\left(\delta_{\mu\rho}\delta_{\nu\sigma}-\delta_{\nu\rho}\delta_{\mu\sigma}\right) \int_0^\infty \mathrm{d}s
            \end{split}
            & &
          \end{flalign}
        \item incoming $\Sigma_E$-gluon-gluon vertex:
          \begin{flalign}
            &
            \begin{gathered}
              \includegraphics{images/E_sigma_vertex_in.pdf}
            \end{gathered}
            & = &
            \begin{split}
              \frac{g^2}{\sqrt{2}} f^{abc}\left(\delta_{\mu\rho}\delta_{\nu\sigma}-\delta_{\nu\rho}\delta_{\mu\sigma}\right)
            \end{split}
            & &
          \end{flalign}
      \end{itemize}
\end{itemize}


\section{Analytical results}\label{app:analytic}

Here, we provide the analytical results for the color coefficients as
far as they are available. For the $\TF\CA$ term of $e_{2,0}$ in
\eqn{eq:econst}, we find
\begin{equation}
  \begin{split}
    -(31.5652\ldots) =&\ \frac{2}{81} \Bigg(-108 \,\text{Li}_3\left(-\frac{1}{3}\right) -792 \,\text{Li}_2\left(\frac{3}{4}\right)+216 \,\text{Li}_2\left(-\frac{1}{3}\right)\\
    &+ 108 \,\text{Li}_{1,2}\left(-3,-\frac{1}{3}\right)-216 \,\text{Li}_{1,2}\left(1,-\frac{1}{3}\right) +216 \,\text{Li}_{2,1}\left(-3,-\frac{1}{3}\right) \\
    &-216 \,\text{Li}_{2,1}\left(1,-\frac{1}{3}\right) +216 \,\text{Li}_{1,1,1}\left(-1,-3,-\frac{1}{3}\right) \\
    &-864 \,\text{Li}_{1,1,1}\left(-1,3,-\frac{1}{3}\right) +216 \,\text{Li}_{1,1,1}\left(1,3,-\frac{1}{3}\right) \\
    &+432 \,\text{Li}_{1,1,1}\left(3,1,-\frac{1}{3}\right) -1134 \zeta (3)+90 \pi ^2-641-54 \ln^3 3 \\
    &+324 \ln^3 2 +324 \ln 2 \ln^2 3+270 \ln^2 3-648 \ln^2 2 \ln 3-2232 \ln^2 2 \\
    &+504 \ln 2 \ln 3-18 \pi ^2 \ln 3+1890\ln 3+108 \pi^2 \ln 2-2148 \ln 2\Bigg) \,,
    \label{eq:e20catf}
  \end{split}
\end{equation}
where
\begin{equation}
  \begin{split}
    \text{Li}_{n_1,\ldots,n_r}(z_1,\ldots,z_r) = \sum_{0<k_1<\cdots<k_r}
    \frac{z_1^{k_1}\cdots z_k^{k_r}}{k_1^{n_1}\cdots k_r^{n_r}}
    %% \label{eq:}
  \end{split}
\end{equation}
are multiple polylogarithms\,\cite{Goncharov:1998kja,Borwein:1999js}.

The $\CF\TF$ coefficient of $C_2$ in \eqn{eq:c2} reads
\begin{equation}
  \begin{split}
    -(3.9226\ldots) = &\ - \frac{1}{18} \bigg(
    48\,\text{Li}_2\left(\frac{1}{9}\right)-2400\,\text{Li}_2
    \left(\frac{1}{3}\right)+672\,\text{Li}_2\left(\frac{3}{4}\right)-131+46
    \pi ^2 \\ &+960 \ln^2 2 -1068 \ln^2 3 +1888 \ln 2-1032\ln 3+624 \ln
    2 \ln 3 \bigg) \,.
    \label{eq:c2cftf}
  \end{split}
\end{equation}

And finally, the $\CF\TF$ coefficient of $s_{2,0}$ in \eqn{eq:sconst}
reads
\begin{equation}
  \begin{split}
    -(15.7975 \dots) =
    &\ \frac{16}{3} \, \text{Li}_2\left(\frac{1}{9}\right) - \frac{800}{3} \, \text{Li}_2\left(\frac{1}{3}\right) + 104 \, \text{Li}_2\left(\frac{3}{4}\right) -\frac{71}{6}+\frac{23 \pi ^2}{9} \\
    &+ 192 \ln^2 2 - \frac{362}{3} \ln^2 3 +\frac{544}{9} \ln 2-\frac{100}{3} \ln 3 + \frac{56}{3} \ln 2 \ln 3 \,.
    \label{eq:s20cftf}
  \end{split}
\end{equation}

%%%% bibliography macros (from source)
\newcommand{\bibentry}[4]{#1, {\it #2}, #3\ifthenelse{\equal{#4}{}}{}{, }#4.}
\newcommand{\journal}[5]{\href{https://dx.doi.org/#5}{\textit{#1} {\bf #2} (#3) #4}}
\newcommand{\arxiv}[2]{\href{https://arXiv.org/abs/#1}{\texttt{arXiv:#1\,[#2]}}}
\newcommand{\arxivhepph}[1]{\href{https://arXiv.org/abs/hep-ph/#1}{\texttt{hep-ph/#1}}}
\newcommand{\arxivhepth}[1]{\href{https://arXiv.org/abs/hep-th/#1}{\texttt{hep-th/#1}}}
\newcommand{\arxivheplat}[1]{\href{https://arXiv.org/abs/hep-lat/#1}{\texttt{hep-lat/#1}}}
\newcommand{\arxivmathph}[1]{\href{https://arXiv.org/abs/math-ph/#1}{\texttt{math-ph/#1}}}
\newcommand{\arxivmath}[1]{\href{https://arXiv.org/abs/math/#1}{\texttt{arXiv:#1 [math]}}}

\begin{thebibliography}{99}
\raggedright
%1
\bibitem{Narayanan:2006rf}
 \bibentry{R. Narayanan and H. Neuberger}%
    {Infinite N phase transitions in continuum Wilson loop operators}%
    {\journal{JHEP}{03}{2006}{064}{10.1088/1126-6708/2006/03/064}}%
    {\arxivhepth{0601210}}%

%2
\bibitem{Luscher:2009eq}
 \bibentry{M. L\"uscher}%
    {Trivializing maps, the Wilson flow and the HMC algorithm}%
    {\journal{Commun. Math. Phys.}{293}{2010}{899}{10.1007/s00220-009-0953-7}}%
    {\arxiv{0907.5491}{hep-lat}}%

%3
\bibitem{Luscher:2010iy}
 \bibentry{M. L\"uscher}%
    {Properties and uses of the Wilson flow in lattice QCD}%
    {\journal{JHEP}{08}{2010}{071}{10.1007/JHEP08(2010)071, 10.1007/JHEP03(2014)092}}%
    {\arxiv{1006.4518}{hep-lat}}%

%4
\bibitem{Luscher:2013vga}
 \bibentry{M. L\"uscher}%
    {Future applications of the Yang-Mills gradient flow in lattice QCD}%
    {\journal{PoS LATTICE2013}{}{2014}{016}{10.22323/1.187.0016}}%
    {\arxiv{1308.5598}{hep-lat}}%

%5
\bibitem{Borsanyi:2012zs}
 \bibentry{S. Bors\'anyi \textit{et al.}}%
    {High-precision scale setting in lattice QCD}%
    {\journal{JHEP}{09}{2012}{010}{10.1007/JHEP09(2012)010}}%
    {\arxiv{1203.4469}{hep-lat}}%

%6
\bibitem{Luscher:2013cpa}
 \bibentry{M. L\"uscher}%
    {Chiral symmetry and the Yang--Mills gradient flow}%
    {\journal{JHEP}{04}{2013}{123}{10.1007/JHEP04(2013)123}}%
    {\arxiv{1302.5246}{hep-lat}}%

%7
\bibitem{Sommer:2014mea}
 \bibentry{R. Sommer}%
    {Scale setting in lattice QCD}%
    {\journal{PoS LATTICE2013}{}{2014}{015}{10.22323/1.187.0015}}%
    {\arxiv{1401.3270}{hep-lat}}%

%8
\bibitem{Ramos:2015dla}
 \bibentry{A. Ramos}%
    {The Yang-Mills gradient flow and renormalization}%
    {\journal{PoS LATTICE2014}{}{2015}{017}{10.22323/1.214.0017}}%
    {\arxiv{1506.00118}{hep-lat}}%

%9
\bibitem{Aoki:2019cca}
 \bibentry{S. Aoki \textit{et al.}}%
    {FLAG Review 2019}%
    {\arxiv{1902.08191}{hep-lat}}{}%

%10
\bibitem{Aoki:2015dla}
 \bibentry{S. Aoki, K. Kikuchi, and T. Onogi}%
    {Geometries from field theories}%
    {\journal{PTEP}{2015}{2015}{101B01}{10.1093/ptep/ptv131}}%
    {\arxiv{1505.00131}{hep-th}}%

%11
\bibitem{Aoki:2016ohw}
 \bibentry{S. Aoki, J. Balog, T. Onogi, and P. Weisz}%
    {Flow equation for the large $N$ scalar model and induced geometries}%
    {\journal{PTEP}{2016}{2016}{083B04}{10.1093/ptep/ptw106}}%
    {\arxiv{1605.02413}{hep-th}}%

%12
\bibitem{Aoki:2017uce}
 \bibentry{S. Aoki and S. Yokoyama}%
    {AdS geometry from CFT on a general conformally flat manifold}%
    {\journal{Nucl. Phys.}{B933}{2018}{262}{10.1016/j.nuclphysb.2018.06.004}}%
    {\arxiv{1709.07281}{hep-th}}%

%13
\bibitem{Aoki:2017bru}
 \bibentry{S. Aoki and S. Yokoyama}%
    {Flow equation, conformal symmetry, and anti-de Sitter geometry}%
    {\journal{PTEP}{2018}{2018}{031B01}{10.1093/ptep/pty013}}%
    {\arxiv{1707.03982}{hep-th}}%

%14
\bibitem{Aoki:2018dmc}
 \bibentry{S. Aoki, J. Balog, and S. Yokoyama}%
    {Holographic computation of quantum corrections to the bulk cosmological constant}%
    {\journal{PTEP}{2019}{2019}{043B06}{10.1093/ptep/ptz026}}%
    {\arxiv{1804.04636}{hep-th}}%

%15
\bibitem{Aoki:2019bfb}
 \bibentry{S. Aoki, S. Yokoyama, and K. Yoshida}%
    {Holographic geometry for nonrelativistic systems emerging from generalized flow equations}%
    {\journal{Phys. Rev.}{D99}{2019}{126002}{10.1103/PhysRevD.99.126002}}%
    {\arxiv{1902.02578}{hep-th}}%

%16
\bibitem{Luscher:2011bx}
 \bibentry{M. L\"uscher and P. Weisz}%
    {Perturbative analysis of the gradient flow in non-abelian gauge theories}%
    {\journal{JHEP}{02}{2011}{051}{10.1007/JHEP02(2011)051}}%
    {\arxiv{1101.0963}{hep-th}}%

%17
\bibitem{Harlander:2016vzb}
 \bibentry{R.V. Harlander and T. Neumann}%
    {The perturbative QCD gradient flow to three loops}%
    {\journal{JHEP}{06}{2016}{161}{10.1007/JHEP06(2016)161}}%
    {\arxiv{1606.03756}{hep-ph}}%

%18
\bibitem{DallaBrida:2017tru}
 \bibentry{M. Dalla Brida and M. L\"uscher}%
    {SMD-based numerical stochastic perturbation theory}%
    {\journal{Eur. Phys. J.}{C77}{2017}{308}{10.1140/epjc/s10052-017-4839-0}}%
    {\arxiv{1703.04396}{hep-lat}}%

%19
\bibitem{Suzuki:2013gza}
 \bibentry{H. Suzuki}%
    {Energy-momentum tensor from the Yang-Mills gradient flow}%
    {\journal{PTEP}{2013}{2013}{083B03}{10.1093/ptep/ptt059, 10.1093/ptep/ptv094}}%
    {\arxiv{1304.0533}{hep-lat}}%

%20
\bibitem{Makino:2014taa}
 \bibentry{H. Makino and H. Suzuki}%
    {Lattice energy-momentum tensor from the Yang-Mills gradient flow---inclusion of fermion fields}%
    {\journal{PTEP}{2014}{2014}{063B02}{10.1093/ptep/ptu070, 10.1093/ptep/ptv095}}%
    {\arxiv{1403.4772}{hep-lat}}%

%21
\bibitem{Harlander:2018zpi}
 \bibentry{R.V. Harlander, Y. Kluth, and F. Lange}%
    {The two-loop energy-momentum tensor within the gradient-flow formalism}%
    {\journal{Eur. Phys. J.}{C78}{2018}{944}{10.1140/epjc/s10052-018-6415-7}}%
    {\arxiv{1808.09837}{hep-lat}}%

%22
\bibitem{Iritani:2018idk}
 \bibentry{T. Iritani, M. Kitazawa, H. Suzuki, and H. Takaura}%
    {Thermodynamics in quenched QCD: energy-momentum tensor with two-loop order coefficients in the gradient-flow formalism}%
    {\journal{PTEP}{2019}{2019}{023B02}{10.1093/ptep/ptz001}}%
    {\arxiv{1812.06444}{hep-lat}}%

%23
\bibitem{Monahan:2016bvm}
 \bibentry{C. Monahan and K. Orginos}%
    {Quasi parton distributions and the gradient flow}%
    {\journal{JHEP}{03}{2017}{116}{10.1007/JHEP03(2017)116}}%
    {\arxiv{1612.01584}{hep-lat}}%

%24
\bibitem{Monahan:2017hpu}
 \bibentry{C. Monahan}%
    {Smeared quasidistributions in perturbation theory}%
    {\journal{Phys. Rev.}{D97}{2018}{054507}{10.1103/PhysRevD.97.054507}}%
    {\arxiv{1710.04607}{hep-lat}}%

%25
\bibitem{Nogueira:1991ex}
 \bibentry{P. Nogueira}%
    {Automatic Feynman graph generation}%
    {\journal{J. Comput. Phys.}{105}{1993}{279}{10.1006/jcph.1993.1074}}%
{}

%26
\bibitem{Nogueira:2006pq}
 \bibentry{P. Nogueira}%
    {Abusing qgraf}%
    {\journal{Nucl. Instrum. Meth.}{A559}{2006}{220}{10.1016/j.nima.2005.11.151}}%
{}

%27
\bibitem{Chetyrkin-priv}
K. Chetyrkin, private communication.

%28
\bibitem{Harlander:1997zb}
 \bibentry{R. Harlander, T. Seidensticker, and M. Steinhauser}%
    {Corrections of $\mathcal{O}(\alpha\alpha_s)$ to the decay of the Z boson
     into bottom quarks}%
    {\journal{Phys. Lett.}{B426}{1998}{125}{10.1016/S0370-2693(98)00220-2}}%
    {\arxivhepph{9712228}}%

%29
\bibitem{Seidensticker:1999bb}
 \bibentry{T. Seidensticker}%
    {Automatic application of successive asymptotic expansions of Feynman diagrams}%
    {\arxivhepph{9905298}}{}%

%30
\bibitem{Vermaseren:2000nd}
 \bibentry{J.A.M. Vermaseren}%
    {New features of FORM}%
    {\arxivmathph{0010025}}{}%

%31
\bibitem{Kuipers:2012rf}
 \bibentry{J. Kuipers, T. Ueda, J.A.M. Vermaseren, and J. Vollinga}%
    {FORM version 4.0}%
    {\journal{Comput. Phys. Commun.}{184}{2013}{1453}{10.1016/j.cpc.2012.12.028}}%
    {\arxiv{1203.6543}{cs.SC}}%

%32
\bibitem{vanRitbergen:1998pn}
 \bibentry{T. van Ritbergen, A.N. Schellekens, and J.A.M. Vermaseren}%
    {Group theory factors for Feynman diagrams}%
    {\journal{Int. J. Mod. Phys.}{A14}{1999}{41}{10.1142/S0217751X99000038}}%
    {\arxivhepph{9802376}}%

%33
\bibitem{Tkachov:1981wb}
 \bibentry{F.V. Tkachov}%
    {A Theorem on Analytical Calculability of Four Loop Renormalization Group Functions}%
    {\journal{Phys. Lett.}{100B}{1981}{65}{10.1016/0370-2693(81)90288-4}}%
{}

%34
\bibitem{Chetyrkin:1981qh}
 \bibentry{K.G. Chetyrkin and F.V. Tkachov}%
    {Integration by Parts: The Algorithm to Calculate beta Functions in 4 Loops}%
    {\journal{Nucl. Phys.}{B192}{1981}{159}{10.1016/0550-3213(81)90199-1}}%
{}

%35
\bibitem{Laporta:2001dd}
 \bibentry{S. Laporta}%
    {High precision calculation of multiloop Feynman integrals by difference equations}%
    {\journal{Int. J. Mod. Phys.}{A15}{2000}{5087}{10.1016/S0217-751X(00)00215-7, 10.1142/S0217751X00002157}}%
    {\arxivhepph{0102033}}%

%36
\bibitem{Mathematica11.3}
  Wolfram Research, Inc., \textit{Mathematica}, Version 11.3, Champaign, IL (2018).

%37
\bibitem{Maierhoefer:2017hyi}
 \bibentry{P. Maierh\"ofer, J. Usovitsch, and P. Uwer}%
    {Kira---A Feynman integral reduction program}%
    {\journal{Comput. Phys. Commun.}{230}{2018}{99}{10.1016/j.cpc.2018.04.012}}%
    {\arxiv{1705.05610}{hep-ph}}%

%38
\bibitem{Maierhofer:2018gpa}
 \bibentry{P. Maierh\"ofer and J. Usovitsch}%
    {Kira 1.2 Release Notes}%
    {\arxiv{1812.01491}{hep-ph}}{}%

%39
\bibitem{Kauers:2008zz}
 \bibentry{M. Kauers}%
    {Fast solvers for dense linear systems}%
    {\journal{Nucl. Phys. Proc. Suppl.}{183}{2008}{245}{10.1016/j.nuclphysbps.2008.09.111}}%
{}

%40
\bibitem{Kant:2013vta}
 \bibentry{P. Kant}%
    {Finding Linear Dependencies in Integration-By-Parts Equations: A Monte Carlo Approach}%
    {\journal{Comput. Phys. Commun.}{185}{2014}{1473}{10.1016/j.cpc.2014.01.017}}%
    {\arxiv{1309.7287}{hep-ph}}%

%41
\bibitem{Peraro:2016wsq}
 \bibentry{T. Peraro}%
    {Scattering amplitudes over finite fields and multivariate functional reconstruction}%
    {\journal{JHEP}{12}{2016}{030}{10.1007/JHEP12(2016)030}}%
    {\arxiv{1608.01902}{hep-ph}}%

%42
\bibitem{Klappert:2019emp}
 \bibentry{J. Klappert and F. Lange}%
    {Reconstructing Rational Functions with $\texttt{FireFly}$}%
    {\arxiv{1904.00009}{cs.SC}}{}%

%43
\bibitem{Zippel:1990}
  \bibentry{R.~Zippel}
  {Interpolating Polynomials from their Values}
  {\journal{Theor. Comp. Sci.}{9}{1990}{375}{10.1016/S0747-7171(08)80018-1}}
  {}

%44
\bibitem{Cuyt:2011}
  \bibentry{A.~Cuyt and W.-s.~Lee}
  {Sparse interpolation of multivariate rational functions}
  {\journal{J. Symb. Comp.}{412}{2011}{1445}{10.1016/j.tcs.2010.11.050}}
  {}

%45
\bibitem{Wang:1981}
  \bibentry{P.S.~Wang}
  {A p-adic Algorithm for Univariate Partial Fractions}
  {\journal{Proc. ACM Symp. Symbolic Algebraic Comp.}{1981}{1981}{212}{10.1145/800206.806398}}
  {}

%46
\bibitem{Monagan:2004}
  \bibentry{M.~Monagan}
  {Maximal Quotient Rational Reconstruction: An Almost Optimal Algorithm for Rational Reconstruction}
  {\journal{Proc. Int. Symp. Symbolic Algebraic Comp.}{2004}{2004}{243}{10.1145/1005285.1005321}}
  {}

%47
\bibitem{Smirnov:2013eza}
 \bibentry{A.V. Smirnov}%
    {FIESTA 3: cluster-parallelizable multiloop numerical calculations in physical regions}%
    {\journal{Comput. Phys. Commun.}{185}{2014}{2090}{10.1016/j.cpc.2014.03.015}}%
    {\arxiv{1312.3186}{hep-ph}}%

%48
\bibitem{Binoth:2003ak}
 \bibentry{T. Binoth and G. Heinrich}%
    {Numerical evaluation of multiloop integrals by sector decomposition}%
    {\journal{Nucl. Phys.}{B680}{2004}{375}{10.1016/j.nuclphysb.2003.12.023}}%
    {\arxivhepph{0305234}}%

%49
\bibitem{Panzer:2014gra}
 \bibentry{E. Panzer}%
    {On hyperlogarithms and Feynman integrals with divergences and many scales}%
    {\journal{JHEP}{03}{2014}{071}{10.1007/JHEP03(2014)071}}%
    {\arxiv{1401.4361}{hep-th}}%

%50
\bibitem{GenzMalik}
	\bibentry{A.~C.~Genz and A.~A.~Malik}
	{An Imbedded Family of Fully Symmetric Numerical Integration Rules}
  {\journal{SIAM J. Numer. Anal.}{20}{1983}{580}{10.1137/0720038}}
	{}

%51
\bibitem{Fousse:2007}
	\bibentry{L. Fousse, G. Hanrot, V. Lef\`{e}vre, P. P{\'e}lissier, and P. Zimmermann}
	{MPFR: A Multiple-precision Binary Floating-point Library with Correct Rounding}
  {\journal{ACM Trans. Math. Softw.}{33}{2007}{}{10.1145/1236463.1236468}}
	{}

%52
\bibitem{Panzer:2014caa}
 \bibentry{E. Panzer}%
    {Algorithms for the symbolic integration of hyperlogarithms with applications to Feynman integrals}%
    {\journal{Comput. Phys. Commun.}{188}{2015}{148}{10.1016/j.cpc.2014.10.019}}%
    {\arxiv{1403.3385}{hep-th}}%

%53
\bibitem{Huber:2005yg}
  \bibentry{T.~Huber and D.~Ma\^itre}
  {\texttt{HypExp}: A Mathematica package for expanding hypergeometric functions around integer-valued parameters}
  {\journal{Comput. Phys. Commun.}{175}{2006}{122}{10.1016/j.cpc.2006.01.007}}
  {\arxivhepph{0507094}}

%54
\bibitem{Huber:2007dx}
  \bibentry{T.~Huber and D.~Ma\^itre}
  {\texttt{HypExp 2}, Expanding Hypergeometric Functions about Half-Integer Parameters}
  {\journal{Comput. Phys. Commun.}{178}{2008}{755}{10.1016/j.cpc.2007.12.008}}
  {\arxiv{0708.2443}{hep-ph}}

%55
\bibitem{Makino:2018rys}
 \bibentry{H. Makino, O. Morikawa, and H. Suzuki}%
    {Gradient flow and the Wilsonian renormalization group flow}%
    {\journal{PTEP}{2018}{2018}{053B02}{10.1093/ptep/pty050}}%
    {\arxiv{1802.07897}{hep-th}}%

%56
\bibitem{Chetyrkin:2000yt}
 \bibentry{K.G. Chetyrkin, J.H. K\"uhn, and M. Steinhauser}%
    {RunDec: A Mathematica package for running and decoupling of the strong coupling and quark masses}%
    {\journal{Comput. Phys. Commun.}{133}{2000}{43}{10.1016/S0010-4655(00)00155-7}}%
    {\arxivhepph{0004189}}%

%57
\bibitem{Herren:2017osy}
 \bibentry{F. Herren and M. Steinhauser}%
    {Version 3 of RunDec and CRunDec}%
    {\journal{Comput. Phys. Commun.}{224}{2018}{333}{10.1016/j.cpc.2017.11.014}}%
    {\arxiv{1703.03751}{hep-ph}}%

%58
\bibitem{Lambrou:2016uet}
 \bibentry{E. Lambrou}%
    {Determining $\alpha_s$ by using the gradient flow in the quenched theory}%
    {\journal{PoS LATTICE2016}{}{2016}{196}{10.22323/1.256.0196}}%
    {}

%59
\bibitem{Ellis:2016jkw}
 \bibentry{J. Ellis}%
    {TikZ-Feynman: Feynman diagrams with TikZ}%
    {\journal{Comput. Phys. Commun.}{210}{2017}{103}{10.1016/j.cpc.2016.08.019}}%
    {\arxiv{1601.05437}{hep-ph}}%

%60
\bibitem{Goncharov:1998kja}
 \bibentry{A.B. Goncharov}%
    {Multiple polylogarithms, cyclotomy and modular complexes}%
    {\journal{Math. Res. Lett.}{5}{1998}{497}{10.4310/MRL.1998.v5.n4.a7}}%
    {\arxiv{1105.2076}{math.AG}}%

%61
\bibitem{Borwein:1999js}
 \bibentry{J.M. Borwein, D.M. Bradley, D.J. Broadhurst, and P. Lisonek}%
    {Special values of multiple polylogarithms}%
    {\journal{Trans. Am. Math. Soc.}{353}{2001}{907}{10.1090/S0002-9947-00-02616-7}}%
    {\arxivmath{9910045}}%
\end{thebibliography}
